% SPDX-License-Identifier: CC-BY-4.0
% Copyright (c) 2026 Martin Schröder <info@swedishembedded.com>

\section{Replicating the result on another repository}\label{sec:replicate}

A result that only holds for the repositories it was developed on is a demonstration, not a
method. This section is the recipe for running the loop on a repository of one's own, with
the status of each step: \emph{available} means a command in the pinned revision does it and
the ledger checks the command line; \emph{planned} means the registration needs it and it
does not exist yet.\cl{PL01} It claims no outcome; the outcome of a replication is the
experiment CE7 of \cref{sec:pending}.

\subsection{Requirements on the repository}

\code{splinter practice} learns a codebase from its own tests: a git work tree or bare
repository, a revision git resolves to a commit, and a test suite that runs offline and is
green and stable on a clean copy of that revision. Splinter recognises a Python project with
a \code{tests} directory (run as \code{python3 -m unittest}) and a Cargo project (run with
\code{cargo test --offline}); any other suite is named with \code{--test-command}, and the files
tasks may change with repeated \code{--file}.\cl{PS01,PS04} The suite is run twice on a clean
copy first, and a red or flaky baseline stops the run before any model is loaded.\cl{PS01}
Isolation is \code{none} unless a container client and a digest-pinned image are configured;
a repository whose tests need isolation for soundness is not suitable on a host without
one.\cl{AU09}

\subsection{Task-family construction}

The protocol's unit is the family (\cref{sec:protocol}). \Cref{tab:families} shows what the
construction yields on the eight authored repositories that ship with Splinter: of several
hundred seeded faults per repository, those that do not build or that no test notices are
dropped, near-duplicates and truncated outputs are refused, and the rest are admitted.\cl{PS02}
For a new repository the steps are, in order:

\begin{enumerate}
\item \textbf{Check the environment} (available): \code{splinter doctor} confirms the state
root is writable with room to spare and the policy's weights fit the card.\cl{PS05}
\item \textbf{Generate and count the admitted tasks without a model} (available): run
\code{splinter practice REPO --rev REV --budget 1s --no-release --max-tasks N
--mutants-per-file M --test-command ...}; the \code{[tasks]} line states the admitted count
and the refusals by reason. This is the funnel of \cref{tab:families}, and it needs no model and
no GPU.\cl{PS03,PS06}
\item \textbf{Reserve the exam} (available): \code{--exam-families} sets families aside by a
stable hash of their names before any attempt, and \code{--exam-seed} is drawn, recorded and
refused if changed on a later run.\cl{PS06}
\item \textbf{Declare partitions and generations} (available for sealed families, planned for
the qualification partition): \code{splinter experiment check} validates a plan whose sealed
families are never trained on; a three-way split with a qualification partition queried once
is the registration's requirement and awaits \code{splinter tasks pick} (planned), which
draws the development holdout and the qualification partition from the task sets by seed and
family.\cl{PS07,PL01}
\item \textbf{Measure the untrained agent} (planned, CE1): the base attempts the development and
qualification partitions with attempts per task fixed from a throughput probe, and
\code{splinter experiences compare CANDIDATE BASELINE} (planned) pairs any two such experience sets
by task and reports the family-clustered sign test and per-family counts.\cl{PL01}
\item \textbf{Train and gate} (available for one round): \code{splinter practice REPO} solves the
non-exam tasks as an agent, keeps the attempts the hidden checks confirm, builds the
action-level SFT dataset, trains a LoRA candidate from the champion with earlier releases
replayed, and runs the release gate over the exam the run reserved; \code{--bf16-base},
\code{--steps}, \code{--rank} and \code{--max-block} bound the training to the card.\cl{PS01,PS06}
\item \textbf{Serve and roll back} (available): a release moves the alias and rewrites the
serving manifest, \code{brain serve --adapter-manifest} follows it, and \code{splinter rollback}
repoints it.\cl{X05,X06,PS08}
\end{enumerate}

\begin{table}[t]
\centering\small
\input{generated/families_table}
\caption{The eight authored repository families: language, tests, and tasks admitted by the
fault-seeding funnel (values read from the family catalogue at the pinned revision and
checked by the ledger).}\label{tab:families}
\end{table}

\subsection{What it costs}

Task generation costs CPU time only: no model is loaded and no GPU is used.\cl{PS03} Training
cost on a 27B model is the token count over the measured rate; at the \mThrNfFourEightOneNineTwoTokensPerS{}~tokens/s of the
probe of \cref{sec:throughput} one million trained tokens take \mThrHoursPerMillionTokens{}~h
of card time.\cl{M38,C12} The remaining terms of the cost of one verified outcome, wall time and
tokens per attempt and attempts per resolved task, are measured by CE1 and are not known
today; they are the reason CE1 precedes any training. A replication should budget by first
running step 2, then CE1 on a small number of families, and only then committing card time.

\subsection{What to report}

For a replication to count it must report, for the new repository, the same quantities as
\cref{sec:pending}: the admission funnel by reason, the baseline resolve rate with its
interval, the paired sign-test counts on the qualification partition, the retention check,
the isolation level and the three repository commits. Counts that disagree with ours are as
informative as ones that agree.\cl{PR13}
